
The network was upgraded with the heat pump, electrode boiler and thermal storage
\emph{after} the measurement record was collected, so validation follows a two-stage
strategy after \citet{Kus2025} and \citet{Maldonado2024}.

\paragraph{Stage 1: direct network validation.}
The pipe model is validated against pre-upgrade monitoring data with the new assets set to
zero capacity. Following the calibration methodology of \citet{Maldonado2024}, the trunk
heat-transfer coefficients $U_p$ are constrained to a defensible range for the pipe class;
this avoids the terminal-pipe inflation that an unconstrained fit to delivered energy would
require, as discussed in Section~\ref{subsec:validation}. Table~\ref{tab:valtargets} lists
the key performance indicators and their gates, fixed \emph{ex ante} from published practice
before the calibration was run.

The use of these gates is deliberate. All are reported in
Section~\ref{subsec:validation}, each against the formulation to which it applies. The
supply-temperature MAE gates are met by the propagating formulation on the validation nodes
(Table~\ref{tab:valtargets}); what the instrumentation cannot support is the \emph{annual,
point} temperature field, whose far-end and corridor comparisons are dominated by a
mixing-valve sensor offset and are therefore reported as diagnostics rather than as gated
metrics. Retaining and reporting these diagnostics is the point rather than an oversight:
they support the claim that the resolution of a model cannot be verified beyond that of its
metering. The conclusions are instead anchored to annual delivered energy, for which the gate
is met; the annual network loss is then \emph{model-implied} (computed from the calibrated
pipe model rather than measured), with delivered-energy agreement fixing the demand scale to
which that loss is referenced, not the loss itself.

\paragraph{Stage 2: necessary-condition verification.}
Dispatch of the new assets is verified against plausibility bounds rather than against
measurements: the heat-pump coefficient of performance must remain within \([2.5,5.5]\),
the thermal-storage state of charge within \([0.05,0.95]\bar{E}_s\), simultaneous charging
and discharging are prohibited, and per-step energy closure is required. These checks are
necessary but not sufficient conditions. A post-upgrade measurement campaign is therefore
the natural follow-up and is identified as a limitation.

\paragraph{Structural sanity checks.}
Two checks confirm internal consistency. The nonlinear reference with its quadratic
constraints deactivated reproduces the baseline objective to within $0.01$\,\%; and the
ordering $\mathrm{cost}(\CP) \le \mathrm{cost}(\Lone)$ holds on every one of the
135 synthetic instances, the expected direction when the copperplate is a loss-free relaxation
of the resolved model on the same network, though not a guaranteed inequality once features
such as CHP export and unit commitment differ between the two. A departure from either check
on this test set would flag a formulation error rather than a finding.
